\subsection{Módulos fielmente planos sobre preesquemas.}
\label{subsection:IV.2.2-fr}

\begin{proposition}[2.2.1]
\label{IV.2.2.1-fr}
Sean $f:X\to Y$ un morfismo de preesquemas y $\mathscr F$ un
$\mathcal O_X$-módulo cuasicoherente. Las condiciones siguientes son
equivalentes:
\begin{enumerate}
\item[\emph{a)}] Para todo cambio de base $g:Y'\to Y$, si se pone
$X'=X\times_Y Y'$, el funtor
$\mathscr G'\leadsto\mathscr F\otimes_Y\mathscr G'$, de la categoría de
$\mathcal O_{Y'}$-módulos cuasicoherentes a la de
$\mathcal O_{X'}$-módulos cuasicoherentes, es exacto y fiel.
\item[\emph{a$'$)}] La condición \emph{a)} se cumple para todos los
morfismos canónicos $g:\operatorname{Spec}(\mathcal O_y)\to Y$ \textup{(\textbf I, 2.4.1)},
donde $y$ recorre $Y$.
\item[\emph{a$''$)}] La condición \emph{a)} se cumple para todas las
inmersiones canónicas $Y'\to Y$, donde $Y'$ recorre el conjunto de los
subesquemas abiertos afines de $Y$.
\item[\emph{b)}] $\mathscr F$ es plano sobre $Y$ y, para todo $y\in Y$,
si se denota por $X_y$ la fibra $f^{-1}(y)$, el
$\mathcal O_{X_y}$-módulo
$\mathscr F_y=\mathscr F\otimes_Y \boldsymbol{k}(y)$ es $\neq0$.
\end{enumerate}
\end{proposition}

Es claro que \emph{a)} implica \emph{a$'$)} y \emph{a$''$)}; la condición
\emph{a$'$)} implica en primer lugar que $\mathscr F$ es plano sobre
$Y$ (2.1.3); por otra parte, implica que para todo $y\in Y$ el funtor
$N\leadsto\mathscr F\otimes_{\mathcal O_y}\widetilde N$ es fiel en la
categoría de los $\mathcal O_y$-módulos; tomando en particular $N=\boldsymbol{k}(y)$,
se obtiene la segunda afirmación de \emph{b)}. Para demostrar que \emph{b)}
implica \emph{a)}, es posible limitarse al caso en que $Y$ es \emph{afín}, pues
la cuestión es local sobre $Y$. Del mismo modo, para demostrar que
\emph{a$''$)} implica \emph{a)}, la cuestión se reduce a probar que, cuando $Y$
es afín, el hecho de que
$\mathscr G\leadsto\mathscr F\otimes_Y\mathscr G$ sea un funtor exacto y fiel
implica la condición \emph{a)}. En otras palabras, la cuestión se reduce a
demostrar la proposición más precisa siguiente:

\begin{proposition}[2.2.2]
\label{IV.2.2.2-fr}
Sean $Y=\operatorname{Spec}(A)$ un esquema afín, $f:X\to Y$ un morfismo de
preesquemas y $\mathscr F$ un $\mathcal O_X$-módulo cuasicoherente. La
condición \emph{a)} de \textup{(2.2.1)} es equivalente a cada una de las siguientes:
\begin{enumerate}
\item[\emph{b$'$)}] $\mathscr F$ es plano sobre $Y$ y, para todo punto \textup{cerrado}
$y$ de $Y$, se tiene $\mathscr F_y\neq0$.
\oldpage[IV-2]{10}
\item[\emph{c)}] El funtor
$\mathscr G\leadsto\mathscr F\otimes_Y\mathscr G$, de la categoría de los
$\mathcal O_Y$-módulos cuasicoherentes a la de los
$\mathcal O_X$-módulos cuasicoherentes, es exacto y fiel.
\end{enumerate}
\end{proposition}

Si se cumple \emph{b$'$)}, existe al menos un $x\in f^{-1}(y)$ tal que
$(\mathscr F_y)_x\neq0$; sea $U=\operatorname{Spec}(B)$ una vecindad abierta
afín de $x$, y sea $\mathscr F|U=\widetilde M$, donde $M$ es un $B$-módulo.
Entonces \emph{b$'$)} implica que $M/\mathfrak j_yM$ es $\neq0$ y, por tanto
(puesto que $M$ es un $A$-módulo plano (2.1.2)), que
$M\otimes_A A_y$ es un $A_y$-módulo fielmente plano
$(\mathbf{0}_{\mathrm I},6.4.5)$. La relación
$(\mathscr F\otimes_Y\mathscr G)\otimes_{\mathcal O_Y}\mathcal O_y=0$ para un punto cerrado
$y$ de $Y$ implica
$(\mathscr F\otimes_{\mathcal O_Y}\mathcal O_y)\otimes_{\mathcal O_y}\mathscr G_y=0$,
y por tanto $\mathscr G_y=0$. Pero, si $\mathscr G_y=0$ para todo punto
cerrado de $Y$, se deduce que $\mathscr G=0$, pues, si
$\mathscr G=\widetilde N$, el anulador de un elemento de $N$ no está contenido
en ningún ideal maximal de $A$, y por tanto es todo $A$. Así, la relación
$\mathscr F\otimes_Y\mathscr G=0$ implica $\mathscr G=0$; en otras palabras,
el funtor $\mathscr G\leadsto\mathscr F\otimes_Y\mathscr G$ es fiel; se sabe,
por otra parte, que este funtor es exacto (2.1.3), lo que demuestra que
\emph{b$'$)} implica \emph{c)}.

Finalmente, para ver que \emph{c)} implica \emph{a)}, es posible limitarse al
caso en que $Y'$ es también afín; como $g:Y'\to Y$ es entonces un morfismo
afín, lo mismo ocurre con la proyección $g':X'\to X$ (\textbf{II}, 1.5.5); además, el
funtor $\mathscr H'\leadsto g'_*(\mathscr H')$ es entonces exacto en la
categoría de los $\mathcal O_{X'}$-módulos cuasicoherentes (\textbf I, 1.6.4), y, si
$g'_*(\mathscr H')=\mathscr H$, se tiene
$\mathscr H'=\widetilde{\mathscr H}$, por lo que el funtor precedente también
es fiel; para ver que
$\mathscr G'\leadsto\mathscr F\otimes_Y\mathscr G'$ es exacto y fiel, basta,
por consiguiente, ver que el funtor
$\mathscr G'\leadsto g'_*(\mathscr F\otimes_Y\mathscr G')$
lo es. Ahora bien, si $f'=f_{(Y')}:X'\to Y'$, se tiene
$g'_*(\mathscr F\otimes_Y\mathscr G')
=g'_*(g'^*(\mathscr F)\otimes_{\mathcal O_{X'}}f'^*(\mathscr G'))$; el hecho de que $g$
sea afín implica que existe un isomorfismo canónico
\begin{equation}
\tag{2.2.2.1}
\label{IV.2.2.2.1-fr}
\mathscr F\otimes_{\mathcal O_X}f^*(g_*(\mathscr G'))
\xrightarrow{\sim}
g'_*(g'^*(\mathscr F)\otimes_{\mathcal O_{X'}}f'^*(\mathscr G')).
\end{equation}
En efecto, se sabe (\textbf{II}, 1.5.2) que existe un isomorfismo canónico
\[
f^*(g_*(\mathscr G'))\xrightarrow{\sim}g'_*(f'^*(\mathscr G')),
\]
y, por otra parte, existe un homomorfismo canónico
$\mathscr F\to g'_*(g'^*(\mathscr F))$ $(\mathbf{0}_{\mathrm I},4.4.3.2)$;
componiendo el homomorfismo
$\mathscr F\otimes_{\mathcal O_X}f^*(g_*(\mathscr G'))
\to g'_*(g'^*(\mathscr F))\otimes_{\mathcal O_X}g'_*(f'^*(\mathscr G'))$
con el homomorfismo canónico $(\mathbf{0}_{\mathrm I},4.2.2.1)$, se obtiene el
homomorfismo (2.2.2.1); la comprobación de que se trata de un isomorfismo es
inmediata reduciéndose al caso en que $X$ es afín. Dicho esto, el funtor
$\mathscr G'\leadsto g_*(\mathscr G')$ es exacto y fiel y, por hipótesis, lo
mismo vale para
$\mathscr G\leadsto\mathscr F\otimes_Y\mathscr G
=\mathscr F\otimes_{\mathcal O_X}f^*(\mathscr G)$; su compuesto es, por tanto,
exacto y fiel, lo que termina de demostrar (2.2.1) y (2.2.2).

\begin{corollary}[2.2.3]
\label{IV.2.2.3-fr}
Sean $X=\operatorname{Spec}(B)$, $Y=\operatorname{Spec}(A)$ dos esquemas
afines, $f:X\to Y$ un morfismo y $\mathscr F=\widetilde M$ un
$\mathcal O_X$-módulo cuasicoherente. Para que $\mathscr F$ cumpla las
condiciones equivalentes de \textup{(2.2.1)} (o \textup{(2.2.2)}), es necesario y suficiente
que el $A$-módulo $M$ sea fielmente plano.
\end{corollary}

En efecto, la condición \emph{c)} de (2.2.2) significa entonces que el funtor
$N\leadsto M\otimes_A N$, de la categoría de los $A$-módulos a la de los
$B$-módulos, es exacto y fiel, y la conclusión se deduce de
$(\mathbf{0}_{\mathrm I},6.4.1)$.

\begin{definition}[2.2.4]
\label{IV.2.2.4-fr}
Cuando se cumplen las condiciones equivalentes de \textup{(2.2.1)}, se dice que el
$\mathcal O_X$-módulo cuasicoherente $\mathscr F$ es fielmente plano
respecto de $f$ (o de $Y$).
\end{definition}

\oldpage[IV-2]{11}
Se observará que esta noción es local sobre $Y$, pero \emph{no sobre} $X$; en
particular, puede ocurrir que $\mathscr F_x=0$ para ciertos $x\in X$; en otras
palabras, $\operatorname{Supp}(\mathscr F)$ no es necesariamente igual a $X$.
Sin embargo, del criterio (2.2.1, \emph{b)}) se deduce inmediatamente que para
todo $y\in Y$ existe al menos un $x\in f^{-1}(y)$ tal que
$(\mathscr F_y)_x=\mathscr F_x\otimes_{\mathcal O_y}\boldsymbol{k}(y)\neq0$, y
\emph{a fortiori} $\mathscr F_x\neq0$; en otros términos:

\begin{corollary}[2.2.5]
\label{IV.2.2.5-fr}
Si $\mathscr F$ es un $\mathcal O_X$-módulo cuasicoherente, fielmente plano
respecto de $f$, se tiene $f(\operatorname{Supp}(\mathscr F))=Y$ y,
\emph{a fortiori}, $f$ es un morfismo sobreyectivo.
\end{corollary}

Este resultado admite una recíproca parcial:

\begin{corollary}[2.2.6]
\label{IV.2.2.6-fr}
Sea $\mathscr F$ un $\mathcal O_X$-módulo cuasicoherente de tipo finito. Para
que $\mathscr F$ sea fielmente plano respecto de $f$, es necesario y suficiente
que $\mathscr F$ sea plano respecto de $f$ y que
$f(\operatorname{Supp}(\mathscr F))=Y$.
\end{corollary}

En efecto, por (\textbf I, 9.1.13 y 3.6.1), se tiene
$\operatorname{Supp}(\mathscr F_y)=f^{-1}(y)\cap\operatorname{Supp}(\mathscr F)$;
por tanto, en este caso el criterio (2.2.1, \emph{b)}) no es sino la condición
del corolario. En particular, el $\mathcal O_X$-módulo $\mathcal O_X$ es
fielmente plano respecto de $f$ si y solo si es plano respecto de $f$ y $f$ es
sobreyectivo; en otras palabras, si y solo si \emph{el morfismo $f$ es
fielmente plano} $(\mathbf{0}_{\mathrm I},6.7.8)$.

Explicitemos las propiedades que intervienen en la definición (2.2.4):

\begin{proposition}[2.2.7]
\label{IV.2.2.7-fr}
Sean $f:X\to Y$ un morfismo de preesquemas y $\mathscr F$ un
$\mathcal O_X$-módulo cuasicoherente fielmente plano respecto de $f$.
Entonces, para que una sucesión
$\mathscr G'\to\mathscr G\to\mathscr G''$ de
$\mathcal O_Y$-módulos cuasicoherentes sea exacta, es necesario y suficiente
que lo sea la sucesión correspondiente
$\mathscr F\otimes_Y\mathscr G'\to\mathscr F\otimes_Y\mathscr G
\to\mathscr F\otimes_Y\mathscr G''$. En particular, para que un homomorfismo
$u:\mathscr G\to\mathscr G'$ de $\mathcal O_Y$-módulos cuasicoherentes sea
inyectivo \textup{(resp.} sobreyectivo, biyectivo, nulo), es necesario y suficiente que
lo sea $1_{\mathscr F}\otimes f^*(u):
\mathscr F\otimes_Y\mathscr G\to\mathscr F\otimes_Y\mathscr G'$. Para que un
$\mathcal O_Y$-módulo cuasicoherente $\mathscr G$ sea nulo, es necesario y
suficiente que $\mathscr F\otimes_Y\mathscr G=0$. Para todo
$\mathcal O_Y$-módulo cuasicoherente $\mathscr G$, la aplicación
$\mathscr G'\leadsto\mathscr F\otimes_Y\mathscr G'$, del conjunto de los
sub-$\mathcal O_Y$-módulos cuasicoherentes de $\mathscr G$ al conjunto de los
sub-$\mathcal O_X$-módulos cuasicoherentes de
$\mathscr F\otimes_Y\mathscr G$, es inyectiva.
\end{proposition}

Para demostrar la última afirmación, es decir, que para dos
sub-$\mathcal O_Y$-módulos cuasicoherentes $\mathscr G'$, $\mathscr G''$ de
$\mathscr G$, la relación
$\mathscr F\otimes_Y\mathscr G'=\mathscr F\otimes_Y\mathscr G''$ implica
$\mathscr G'=\mathscr G''$, es posible (reemplazando $\mathscr G''$ por
$\mathscr G'+\mathscr G''$) reducirse al caso en que
$\mathscr G'\subset\mathscr G''$, y basta entonces aplicar la segunda
afirmación del enunciado a la inyección $u:\mathscr G'\to\mathscr G''$.

\begin{corollary}[2.2.8]
\label{IV.2.2.8-fr}
Sea $f:X\to Y$ un morfismo fielmente plano. Para todo
$\mathcal O_Y$-módulo cuasicoherente $\mathscr G$, la aplicación canónica
\begin{equation}
\tag{2.2.8.1}
\label{IV.2.2.8.1-fr}
\Gamma(Y,\mathscr G)\longrightarrow\Gamma(X,f^*(\mathscr G))
\end{equation}
es inyectiva.
\end{corollary}

En efecto, $\Gamma(Y,\mathscr G)$ se identifica canónicamente con
$\operatorname{Hom}_{\mathcal O_Y}(\mathcal O_Y,\mathscr G)$
$(\mathbf{0}_{\mathrm I},5.1.1)$ y asimismo $\Gamma(X,f^*(\mathscr G))$ con
$\operatorname{Hom}_{\mathcal O_X}(\mathcal O_X,f^*(\mathscr G))$. En virtud
de (2.2.1) y (2.2.4), la hipótesis implica que el funtor
$\mathscr G\leadsto f^*(\mathscr G)$ es exacto y fiel en la categoría de los
$\mathcal O_Y$-módulos cuasicoherentes; por tanto, un homomorfismo
$u:\mathcal O_Y\to\mathscr G$ es nulo si y solo si
$f^*(u):f^*(\mathcal O_Y)=\mathcal O_X\to f^*(\mathscr G)$ es nulo.
